\documentclass[11pt]{article}
\usepackage{amsmath, amssymb, amsthm}

\title{A6 Meta-Analysis: System Interaction and Global Geometry}
\author{Borealis S. Hedling}
\date{October 2026}

\begin{document}
\maketitle

\section*{Abstract}
This document provides a meta-analysis of the A6 dual-algebra system. It explains
how geometric definitions, invariants, operators, unified kernel geometry,
runtime diagnostics, and the safety envelope interact to form a stable,
invariant-preserving computational geometry. This analysis prepares the system
for A7 divergence membranes by demonstrating that A6 is structurally complete,
runtime-safe, and invariant-stable.

\section{Overview of A6 Architecture}
A6 consists of ten Lean modules and three physics layers:
\begin{itemize}
  \item geometric definitions and invariants,
  \item operator algebra and unified semantics,
  \item unified kernel, diagnostics, and safety envelope.
\end{itemize}

These layers form a dependency chain:
\[
\text{Geometry} \to \text{Invariants} \to \text{Operators} \to
\text{Semantics} \to \text{Kernel} \to \text{Diagnostics} \to \text{Safety}.
\]

\section{Invariant Propagation}
Invariants propagate through the operator stack:
\[
\text{Spectral} \to \text{Modal} \to \text{Holonomy} \to
\text{Collapse} \to \text{Stability}.
\]

Each invariant restricts the next:
\begin{itemize}
  \item spectral continuity restricts modal adjacency,
  \item modal adjacency restricts holonomy drift,
  \item holonomy drift restricts collapse geometry,
  \item collapse geometry restricts stability envelopes.
\end{itemize}

This chain ensures that operator composition remains bounded.

\section{Operator Interaction}
Operators interact through unified semantics:
\[
O_{\text{unified}} = 
(O_{\text{spec}}, O_{\text{modal}}, O_{\text{holo}},
 O_{\text{collapse}}, O_{\text{stab}}).
\]

Operator composition must preserve:
\begin{itemize}
  \item kernel membership,
  \item invariant structure,
  \item diagnostic validity,
  \item safety-envelope bounds.
\end{itemize}

\section{Unified Kernel Interaction}
The unified kernel $K$ is the intersection of all operator kernels:
\[
K = K_{\text{spec}} \cap K_{\text{modal}} \cap K_{\text{holo}} \cap
    K_{\text{collapse}} \cap K_{\text{stab}}.
\]

The kernel acts as a firewall:
\[
p \in K \implies O(p) \in K.
\]

This ensures that operator execution cannot escape the invariant-preserving
region of the geometry.

\section{Collapse-Surface Interaction}
Collapse surfaces $\Sigma$ intersect the kernel transversally:
\[
T_p K + T_p \Sigma = T_p X.
\]

This ensures:
\begin{itemize}
  \item collapse detection is well-defined,
  \item collapse surfaces cannot degenerate inside the kernel,
  \item collapse events cannot propagate uncontrollably.
\end{itemize}

\section{Stability-Envelope Interaction}
The stability envelope $E$ bounds perturbations:
\[
|E(p) - E(q)| < \epsilon_{\text{stab}}.
\]

This ensures:
\begin{itemize}
  \item operator perturbations remain bounded,
  \item collapse surfaces remain stable,
  \item holonomy drift cannot accumulate,
  \item spectral mismatch cannot blow out.
\end{itemize}

\section{Runtime Diagnostics}
Diagnostics monitor:
\begin{itemize}
  \item spectral mismatch,
  \item modal discontinuity,
  \item holonomy drift,
  \item collapse proximity,
  \item stability-envelope deviation,
  \item kernel firewall integrity.
\end{itemize}

Diagnostics run in parallel with operator execution:
\[
\text{RuntimeState}(p) =
(O(p), D_{\text{spec}}, D_{\text{modal}}, D_{\text{holo}},
 D_{\text{collapse}}, D_{\text{stab}}, D_{\text{kernel}}).
\]

Warnings are raised when any diagnostic exceeds its threshold.

\section{Safety Envelope}
The safety envelope $S$ wraps the entire system:
\[
S = \{ p \mid E(p) \leq \epsilon_S \}.
\]

It ensures:
\begin{itemize}
  \item operators remain bounded,
  \item diagnostics remain valid,
  \item kernel remains inside the safety region,
  \item collapse surfaces cannot occur inside $S$,
  \item holonomy drift cannot destabilize the geometry.
\end{itemize}

\section{A6 Completeness}
A6 is complete when:
\[
\text{Kernel} \land \text{Diagnostics} \land \text{Safety}
\]
are all satisfied.

This ensures that A6 is:
\begin{itemize}
  \item invariant-preserving,
  \item operator-stable,
  \item collapse-safe,
  \item holonomy-bounded,
  \item spectrally-continuous,
  \item modal-adjacent,
  \item runtime-safe.
\end{itemize}

\section{Preparation for A7 Divergence Membrane}
A6 completeness is required before studying A7 divergence membranes. A7
introduces:
\begin{itemize}
  \item divergence geometry,
  \item divergence invariants,
  \item divergence operators,
  \item divergence collapse surfaces,
  \item divergence stability envelopes,
  \item divergence kernel.
\end{itemize}

A6 ensures that the system is stable enough to safely explore A7.

\section*{Conclusion}
This meta-analysis demonstrates that A6 is structurally complete, invariant-
preserving, operator-stable, kernel-safe, diagnostic-valid, and safety-bounded.
It provides the global understanding required for future researchers and
prepares the system for A7 divergence membranes.

\end{document}
